% Part: methods
% Chapter: induction
% Section: introduction

\documentclass[../../../include/open-logic-section]{subfiles}

\begin{document}

\olfileid[id]{mth}{ind}{int}

\olsection{Pendahuluan}

Induksi adalah teknik pembuktian penting yang, dalam berbagai bentuk,
digunakan di hampir semua bidang logika, ilmu komputer teoretis, dan
matematika. Teknik ini diperlukan untuk membuktikan banyak hasil dalam
logika.

Induksi sering dipertentangkan dengan deduksi dan dicirikan sebagai
inferensi dari yang khusus menuju yang umum. Misalnya, jika kita mengamati
banyak zamrud hijau dan tidak menemukan apa pun yang akan kita sebut zamrud
tetapi tidak berwarna hijau, kita mungkin menyimpulkan bahwa semua zamrud
berwarna hijau. Ini merupakan inferensi induktif karena bergerak dari banyak
kasus khusus (zamrud ini hijau, zamrud itu hijau, dan seterusnya) menuju
suatu klaim umum (semua zamrud hijau). Induksi \emph{matematis} juga
merupakan inferensi yang menghasilkan suatu klaim umum, tetapi jenisnya
sangat berbeda dari ``induksi sederhana'' ini.

Secara sangat kasar, bukti induktif dalam matematika menyimpulkan bahwa
semua objek matematis dari jenis tertentu memiliki suatu sifat. Dalam kasus
paling sederhana, objek matematis yang dibahas oleh bukti induktif adalah
bilangan asli. Dalam kasus itu, bukti induktif digunakan untuk menetapkan
bahwa semua bilangan asli memiliki suatu sifat, dengan menunjukkan bahwa
\begin{enumerate}
    \item $0$ memiliki sifat tersebut, dan
    \item setiap kali suatu bilangan~$k$ memiliki sifat tersebut, $k+1$
      juga memilikinya.
\end{enumerate}
Induksi pada bilangan asli juga sering dapat digunakan untuk membuktikan
klaim umum mengenai objek matematis yang dapat diberi bilangan. Misalnya,
setiap himpunan hingga memiliki sejumlah hingga~$n$ anggota; jika dengan
induksi kita dapat menunjukkan bahwa setiap bilangan~$n$ memiliki sifat
``semua himpunan hingga berukuran~$n$ adalah \dots'', kita telah menunjukkan
sesuatu mengenai semua himpunan hingga.

Induksi juga dapat diperumum ke objek matematis yang \emph{didefinisikan
secara induktif}. Misalnya, ungkapan suatu bahasa formal, seperti ungkapan
dalam logika orde pertama, didefinisikan secara induktif. \emph{Induksi
struktural} merupakan cara untuk membuktikan hasil mengenai semua ungkapan
tersebut. Induksi struktural khususnya sangat berguna---dan digunakan secara
luas---dalam logika.

\end{document}
